\section{What the regulations actually say}
\label{sec:regulatory-asymmetry}

How far the regulatory texts support the distinction depends on the jurisdiction. United States policy builds the \mention{from}/\mention{about} contrast into its definition. Under the Common Rule, a human subject is a living individual \enquote{about whom} an investigator obtains data through intervention or interaction, or whose identifiable private information is used \citep{commonrule2018}. A rater's score can be obtained \emph{from} a person without being \emph{about} that person.

Canadian policy offers no such phrase. Canada's Tri-Council Policy Statement (TCPS 2) defines participants as individuals whose data or \enquote{responses to interventions, stimuli or questions by the researcher, are relevant to answering the research question(s)} \citep{tcps2_2022}. A rater's judgment answers a question and is plainly relevant to the research question, so on its face the definition takes in the evaluator. The distinction can't rest on the bare fact that a judgment was requested.

It needn't. TCPS 2 itself allows interaction with people who aren't the focus of the research: it points to \enquote{authorized personnel} who release role-based information \enquote{about organizations, policies, procedures, professional practices or statistical reports}, and says \enquote{such individuals are not considered participants} \citep{tcps2_2022}. The line is drawn by focus of inquiry, not by an \mention{about} clause.

The Panel's interpretation carries the point to where it bites: research relying only on information that staff \enquote{normally provide as part of their work duties} needs no review, because \enquote{the information is the focus of the research, not the views of the staff member} \citep{pre_scope}. An expert asked whether an item meets a stated criterion is supplying information of just that kind: the focus is the item, not the assessor.

The parallel doesn't require an external evaluator to be staff, or the judgment to be an employment duty. What carries over is the focus-of-inquiry principle: information supplied in a role can be the research's focus without the supplier becoming its object.

The same interpretation marks the boundary the argument has to respect. If researchers also ask those people \enquote{to provide personal opinions outside the scope of their job roles}, review is required \citep{pre_scope}. The evaluator role holds only while the judgment is read as a within-competence assessment of the object. Once a judgment is taken as the person's own reaction, idiolect, or attitude, the focus has moved to the person, and the judge becomes a participant (\autoref{sec:boundary-cases}).

This locates the point about method correctly. That \enquote{choice of methodology and/or intent or ability to publish findings} don't determine whether an activity is research \citep{tcps2_2022} blocks the inference from systematicity to participant status, but it doesn't by itself separate participant from evaluator. The focus-of-inquiry test does that work. The judge evaluates the item; the item doesn't evaluate the judge.
